% MATH3701 Supplementary Proofs — shared preamble
% % MATH3701 Supplementary Proofs — shared preamble
% % MATH3701 Supplementary Proofs — shared preamble
% \input{proof_preamble} at the top of each proof file

\documentclass[11pt,a4paper]{article}
\usepackage{amsmath,amssymb,amsthm}
\usepackage{mathtools}
\usepackage{bm}
\usepackage{geometry}
\geometry{margin=2.5cm}
\usepackage{hyperref}
\hypersetup{colorlinks=true, linkcolor=blue, urlcolor=blue}

% ---- Math shortcuts (matching slides/preamble.tex) ----
\newcommand{\E}{\mathbb{E}}
\newcommand{\Var}{\mathrm{Var}}
\newcommand{\Cov}{\mathrm{Cov}}
\newcommand{\Prob}{\mathbb{P}}
\newcommand{\R}{\mathbb{R}}
\newcommand{\N}{\mathcal{N}}
\newcommand{\bX}{\mathbf{X}}
\newcommand{\bx}{\mathbf{x}}
\newcommand{\by}{\mathbf{y}}
\newcommand{\bY}{\mathbf{Y}}
\newcommand{\bb}{\mathbf{b}}
\newcommand{\ba}{\mathbf{a}}
\newcommand{\bbeta}{\boldsymbol{\beta}}
\newcommand{\beps}{\boldsymbol{\epsilon}}
\newcommand{\bmu}{\boldsymbol{\mu}}
\newcommand{\btheta}{\boldsymbol{\theta}}
\newcommand{\bU}{\mathbf{U}}
\newcommand{\bJ}{\mathbf{J}}
\newcommand{\bI}{\mathbf{I}}
\newcommand{\bH}{\mathbf{H}}
\newcommand{\bK}{\mathbf{K}}
\newcommand{\bL}{\mathbf{L}}
\newcommand{\bM}{\mathbf{M}}
\newcommand{\bS}{\mathbf{S}}
\newcommand{\bW}{\mathbf{W}}
\newcommand{\bz}{\mathbf{z}}
\newcommand{\boldeta}{\boldsymbol{\eta}}
\newcommand{\logit}{\mathrm{logit}}
\newcommand{\iid}{\stackrel{\mathrm{iid}}{\sim}}
\newcommand{\tr}{\mathrm{tr}}

% ---- Theorem environments ----
\theoremstyle{plain}
\newtheorem{proposition}{Proposition}
\newtheorem{lemma}{Lemma}
\newtheorem{theorem}{Theorem}
\theoremstyle{definition}
\newtheorem{remark}{Remark}

% ---- Module branding ----
\newcommand{\moduleheader}{%
  \noindent\textbf{MATH3701 Statistical Modelling}\\
  School of Mathematics, University of Leeds\\[6pt]
}
 at the top of each proof file

\documentclass[11pt,a4paper]{article}
\usepackage{amsmath,amssymb,amsthm}
\usepackage{mathtools}
\usepackage{bm}
\usepackage{geometry}
\geometry{margin=2.5cm}
\usepackage{hyperref}
\hypersetup{colorlinks=true, linkcolor=blue, urlcolor=blue}

% ---- Math shortcuts (matching slides/preamble.tex) ----
\newcommand{\E}{\mathbb{E}}
\newcommand{\Var}{\mathrm{Var}}
\newcommand{\Cov}{\mathrm{Cov}}
\newcommand{\Prob}{\mathbb{P}}
\newcommand{\R}{\mathbb{R}}
\newcommand{\N}{\mathcal{N}}
\newcommand{\bX}{\mathbf{X}}
\newcommand{\bx}{\mathbf{x}}
\newcommand{\by}{\mathbf{y}}
\newcommand{\bY}{\mathbf{Y}}
\newcommand{\bb}{\mathbf{b}}
\newcommand{\ba}{\mathbf{a}}
\newcommand{\bbeta}{\boldsymbol{\beta}}
\newcommand{\beps}{\boldsymbol{\epsilon}}
\newcommand{\bmu}{\boldsymbol{\mu}}
\newcommand{\btheta}{\boldsymbol{\theta}}
\newcommand{\bU}{\mathbf{U}}
\newcommand{\bJ}{\mathbf{J}}
\newcommand{\bI}{\mathbf{I}}
\newcommand{\bH}{\mathbf{H}}
\newcommand{\bK}{\mathbf{K}}
\newcommand{\bL}{\mathbf{L}}
\newcommand{\bM}{\mathbf{M}}
\newcommand{\bS}{\mathbf{S}}
\newcommand{\bW}{\mathbf{W}}
\newcommand{\bz}{\mathbf{z}}
\newcommand{\boldeta}{\boldsymbol{\eta}}
\newcommand{\logit}{\mathrm{logit}}
\newcommand{\iid}{\stackrel{\mathrm{iid}}{\sim}}
\newcommand{\tr}{\mathrm{tr}}

% ---- Theorem environments ----
\theoremstyle{plain}
\newtheorem{proposition}{Proposition}
\newtheorem{lemma}{Lemma}
\newtheorem{theorem}{Theorem}
\theoremstyle{definition}
\newtheorem{remark}{Remark}

% ---- Module branding ----
\newcommand{\moduleheader}{%
  \noindent\textbf{MATH3701 Statistical Modelling}\\
  School of Mathematics, University of Leeds\\[6pt]
}
 at the top of each proof file

\documentclass[11pt,a4paper]{article}
\usepackage{amsmath,amssymb,amsthm}
\usepackage{mathtools}
\usepackage{bm}
\usepackage{geometry}
\geometry{margin=2.5cm}
\usepackage{hyperref}
\hypersetup{colorlinks=true, linkcolor=blue, urlcolor=blue}

% ---- Math shortcuts (matching slides/preamble.tex) ----
\newcommand{\E}{\mathbb{E}}
\newcommand{\Var}{\mathrm{Var}}
\newcommand{\Cov}{\mathrm{Cov}}
\newcommand{\Prob}{\mathbb{P}}
\newcommand{\R}{\mathbb{R}}
\newcommand{\N}{\mathcal{N}}
\newcommand{\bX}{\mathbf{X}}
\newcommand{\bx}{\mathbf{x}}
\newcommand{\by}{\mathbf{y}}
\newcommand{\bY}{\mathbf{Y}}
\newcommand{\bb}{\mathbf{b}}
\newcommand{\ba}{\mathbf{a}}
\newcommand{\bbeta}{\boldsymbol{\beta}}
\newcommand{\beps}{\boldsymbol{\epsilon}}
\newcommand{\bmu}{\boldsymbol{\mu}}
\newcommand{\btheta}{\boldsymbol{\theta}}
\newcommand{\bU}{\mathbf{U}}
\newcommand{\bJ}{\mathbf{J}}
\newcommand{\bI}{\mathbf{I}}
\newcommand{\bH}{\mathbf{H}}
\newcommand{\bK}{\mathbf{K}}
\newcommand{\bL}{\mathbf{L}}
\newcommand{\bM}{\mathbf{M}}
\newcommand{\bS}{\mathbf{S}}
\newcommand{\bW}{\mathbf{W}}
\newcommand{\bz}{\mathbf{z}}
\newcommand{\boldeta}{\boldsymbol{\eta}}
\newcommand{\logit}{\mathrm{logit}}
\newcommand{\iid}{\stackrel{\mathrm{iid}}{\sim}}
\newcommand{\tr}{\mathrm{tr}}

% ---- Theorem environments ----
\theoremstyle{plain}
\newtheorem{proposition}{Proposition}
\newtheorem{lemma}{Lemma}
\newtheorem{theorem}{Theorem}
\theoremstyle{definition}
\newtheorem{remark}{Remark}

% ---- Module branding ----
\newcommand{\moduleheader}{%
  \noindent\textbf{MATH3701 Statistical Modelling}\\
  School of Mathematics, University of Leeds\\[6pt]
}


\begin{document}

\moduleheader
\textbf{Supplementary Proof: The Cox Partial Likelihood}\\[6pt]
\textit{This note justifies why the partial likelihood provides a valid basis for inference on $\bbeta$ in the Cox model, without requiring specification of the baseline hazard $h_0(t)$.}

\bigskip

\section*{Setup}

The Cox proportional hazards model is
\[
h(t \mid \bx_i) = h_0(t) \exp(\bbeta^T \bx_i), \qquad i = 1, \ldots, n.
\]
We observe $(t_i, \delta_i, \bx_i)$ where $t_i$ is the observed time, $\delta_i$ the event indicator, and $\bx_i \in \mathbb{R}^p$ the covariate vector.

Let $t_{(1)} < t_{(2)} < \cdots < t_{(r)}$ be the ordered distinct event times (assuming no ties for simplicity), and let $\mathcal{R}(t_{(j)}) = \{i : t_i \geq t_{(j)}\}$ be the risk set just before $t_{(j)}$.

\section*{Derivation of the partial likelihood}

Cox (1972) proposed to condition on the event times $t_{(1)}, \ldots, t_{(r)}$ and use only the information about which individual experienced the event at each time.

\begin{lemma}
Under the Cox model, the conditional probability that individual $i_{(j)}$ is the one who dies at time $t_{(j)}$, given that exactly one death occurs in $\mathcal{R}(t_{(j)})$ at that time, is
\[
\Prob\bigl(T_{i_{(j)}} = t_{(j)} \mid \text{one death in } \mathcal{R}(t_{(j)})\bigr)
= \frac{\exp(\bbeta^T \bx_{i_{(j)}})}{\displaystyle\sum_{\ell \in \mathcal{R}(t_{(j)})} \exp(\bbeta^T \bx_\ell)}.
\]
\end{lemma}

\begin{proof}
For a small interval $[t, t + \delta t)$, the probability that individual $\ell$ dies in this interval, given they are at risk, is approximately $h(t \mid \bx_\ell)\,\delta t = h_0(t)\exp(\bbeta^T\bx_\ell)\,\delta t$.

The probability that individual $i_{(j)}$ dies at $t_{(j)}$ while all others in the risk set survive past $t_{(j)}$ is proportional to $h_0(t_{(j)})\exp(\bbeta^T\bx_{i_{(j)}})$. The probability that \emph{some} individual in $\mathcal{R}(t_{(j)})$ dies at $t_{(j)}$ is proportional to $h_0(t_{(j)})\sum_{\ell \in \mathcal{R}} \exp(\bbeta^T\bx_\ell)$.

Dividing:
\[
\Prob(i_{(j)} \text{ dies} \mid \text{one death in } \mathcal{R})
= \frac{h_0(t_{(j)})\exp(\bbeta^T\bx_{i_{(j)}})}{h_0(t_{(j)})\sum_{\ell \in \mathcal{R}(t_{(j)})} \exp(\bbeta^T\bx_\ell)}
= \frac{\exp(\bbeta^T\bx_{i_{(j)}})}{\sum_{\ell \in \mathcal{R}(t_{(j)})} \exp(\bbeta^T\bx_\ell)}.
\]
The baseline hazard $h_0(t_{(j)})$ cancels completely.
\end{proof}

\begin{remark}
This cancellation is the key insight: by conditioning on the observed event times, we eliminate the nuisance function $h_0(t)$ entirely. The remaining expression depends only on $\bbeta$.
\end{remark}

\section*{The partial likelihood}

Multiplying the conditional probabilities across all $r$ event times, the \textbf{partial likelihood} is

\[
L(\bbeta) = \prod_{j=1}^r \frac{\exp(\bbeta^T\bx_{(j)})}{\displaystyle\sum_{\ell \in \mathcal{R}(t_{(j)})} \exp(\bbeta^T\bx_\ell)},
\]

or equivalently, using the event indicators $\delta_i$:

\[
L(\bbeta) = \prod_{i=1}^n \left[\frac{\exp(\bbeta^T\bx_i)}{\sum_{\ell \in \mathcal{R}(t_i)} \exp(\bbeta^T\bx_\ell)}\right]^{\!\delta_i}.
\]

The corresponding partial log-likelihood is

\[
\ell(\bbeta) = \sum_{i=1}^n \delta_i \left\{\bbeta^T\bx_i - \log \sum_{\ell \in \mathcal{R}(t_i)} \exp(\bbeta^T\bx_\ell)\right\}.
\]

\section*{Score equations}

Differentiating $\ell(\bbeta)$ with respect to $\beta_k$:

\[
\frac{\partial \ell}{\partial \beta_k} = \sum_{i=1}^n \delta_i \left\{x_{ik} - \frac{\sum_{\ell \in \mathcal{R}(t_i)} x_{\ell k} \exp(\bbeta^T\bx_\ell)}{\sum_{\ell \in \mathcal{R}(t_i)} \exp(\bbeta^T\bx_\ell)}\right\} = 0.
\]

These $p$ equations are solved numerically (Newton--Raphson) to give the maximum partial likelihood estimator (MPLE) $\hat\bbeta$.

\section*{Asymptotic properties}

\begin{theorem}[Cox 1975; Tsiatis 1981]
Under standard regularity conditions, the MPLE $\hat\bbeta$ satisfies:
\begin{enumerate}
\item \textbf{Consistency:} $\hat\bbeta \xrightarrow{\Prob} \bbeta$ as $n \to \infty$.
\item \textbf{Asymptotic normality:}
\[
\sqrt{n}(\hat\bbeta - \bbeta) \xrightarrow{d} \mathcal{N}_p\!\left(\mathbf{0},\, \mathcal{I}(\bbeta)^{-1}\right),
\]
where $\mathcal{I}(\bbeta) = -n^{-1}\,\partial^2\ell/\partial\bbeta\partial\bbeta^T$ is the (per-observation) observed information matrix evaluated at $\hat\bbeta$.
\item \textbf{Efficiency:} Among all estimators based on the rank ordering of event times, $\hat\bbeta$ is asymptotically efficient.
\end{enumerate}
\end{theorem}

The asymptotic normality justifies the Wald confidence intervals reported by \texttt{coxph()} in R:
\[
\hat\beta_k \pm z_{1-\alpha/2} \cdot \widehat{\mathrm{sd}}(\hat\beta_k),
\]
where $\widehat{\mathrm{sd}}(\hat\beta_k) = \sqrt{[\mathcal{I}(\hat\bbeta)^{-1}]_{kk}}$.

\section*{Partial likelihood ratio test}

For nested models $M_0 \subset M_1$ with $q - p$ additional parameters in $M_1$, the partial likelihood ratio test statistic
\[
\Lambda = -2(\hat\ell_0 - \hat\ell_1) \;\dot\sim\; \chi^2_{q-p}
\]
has the same asymptotic $\chi^2$ distribution as for standard MLEs, justifying the use of this test in \texttt{coxph()} comparisons.

\section*{Tied event times: Breslow's approximation}

With $d_j$ tied deaths at $t_{(j)}$ with covariate sum $\mathbf{s}_{(j)} = \sum_{k : t_k = t_{(j)}, \delta_k = 1} \bx_k$, Breslow's approximation replaces the exact combinatorial expression with

\[
L(\bbeta) \approx \prod_{j=1}^r \frac{\exp(\bbeta^T\mathbf{s}_{(j)})}{\left(\sum_{\ell \in \mathcal{R}(t_{(j)})} \exp(\bbeta^T\bx_\ell)\right)^{d_j}}.
\]

This approximation is accurate when ties are rare (the typical situation in continuous survival data).

\section*{References}

Cox, D.R.\ (1972). Regression models and life-tables (with discussion). \textit{Journal of the Royal Statistical Society, Series B}, \textbf{34}, 187--220.

Cox, D.R.\ (1975). Partial likelihood. \textit{Biometrika}, \textbf{62}, 269--276.

Tsiatis, A.A.\ (1981). A large sample study of Cox's regression model. \textit{The Annals of Statistics}, \textbf{9}, 93--108.

\end{document}
